\scanpage{54}
\mspage{53}
2 - 5 - \hfill n° 392
\[
\begin{aligned}
(2)\qquad &\underline{ab}=\prod\underline p^{m(\underline p)+n(\underline p)}\\
(3)\qquad &\underline a:\underline b=\underline{ab}^{-1}=\prod\underline p^{m(\underline p)-n(\underline p)}\\
(4)\qquad &\underline a+\underline b=\prod\underline p^{\inf(m(\underline p),n(\underline p))}\\
(5)\qquad &\underline a\cap\underline b=\prod\underline p^{\sup(m(\underline p),n(\underline p))}.
\end{aligned}
\]
En effet, la relation (2) est évidente ; la relation (3) en résulte, l'égalité $\underline a:\underline b=\underline{ab}^{-1}$ découlant de la formule $\operatorname{div}(\underline a:\underline b)=\operatorname{div}\underline a-\operatorname{div}\underline b$, (§ 1, n° 2, cor. du th. 1) ; les formules (4) et (5) résultent de la prop. 2, § 1, n° 1.

Ces résultats s'appliquent notamment à la clôture intégrale de $Z$ dans une extension de degré fini de $Q$.

Lorsque $A$ est principal, les résultats ci-dessus redonnent ceux d'\uline{Alg.}, chap. VII, § 1, n° 3.

4. \uline{Théorème d'approximation dans les anneaux de Dedekind}.

Dans les anneaux de Dedekind, on a un ``théorème d'approximation'' qui améliore à la fois le th. 1 du chap. VI, § 7, n° 2 et la prop. 9 du § 1, n° 5 :

PROPOSITION 2.- \uline{Soient} $A$ \uline{un anneau de Dedekind}, $K$ \uline{son corps des fractions}, \uline{et} $P$ \uline{l'ensemble des idéaux premiers non nuls de} $A$ ; \uline{pour} $\underline p\in P$, \uline{notons} $v_{\underline p}$ \uline{la valuation essentielle correspondante de} $A$. \uline{Soient} $\underline p_1,\ldots,\underline p_q$ \uline{des éléments deux à deux distincts de} $P$, $n_1,\ldots,n_q$ \uline{des entiers rationnels}, \uline{et} $x_1,\ldots,x_q$ \uline{des éléments de} $K$. \uline{Il existe alors} $x\in K$ \uline{tel que} $v_{\underline p_i}(x-x_i)\ge n_i$ \uline{pour} $1\le i\le q$, \uline{et que} $v_{\underline p}(x)\ge0$ \uline{pour tout} $\underline p\in P$ \uline{distinct des} $\underline p_i$.

En remplaçant au besoin les $n_i$ par des entiers qui leurs sont supérieurs, on peut les supposer tous positifs. Examinons d'abord le cas où les $x_i$ sont dans $A$ ; tout revient évidemment à trouver un $x\in A$ vérifiant les congruences
\[
x\equiv x_i\ (\mathrm{mod.}\ \underline p^{n_i})
\]
et l'existence de $x$ résulte alors du chap. II, § 1, n° 2, prop. 5.
